\begin{exercise}\label{exr:vi26-01}Show that a quotient of a finitely generated $A$-algebra is finitely generated over $A$.\end{exercise}

\begin{exercise}\label{exr:vi26-02}Show $A_f$ is a finite-type $A$-algebra.\end{exercise}

\begin{exercise}\label{exr:vi26-03}If $B$ is finite type over $A$ and $C$ finite type over $B$, show $C$ finite type over $A$.\end{exercise}

\begin{exercise}\label{exr:vi26-04}Show $A[x_1,\dots,x_n]/I$ is finite type over $A$ for arbitrary $I$.\end{exercise}

\begin{exercise}\label{exr:vi26-05}Show that finite type need not imply finite presentation over a non-Noetherian base.\end{exercise}

\begin{exercise}\label{exr:vi26-06}Prove a finite-type morphism is quasi-compact.\end{exercise}

\begin{exercise}\label{exr:vi26-07}Show an affine morphism locally of finite type is finite type.\end{exercise}

\begin{exercise}\label{exr:vi26-08}Show $\mathbb A^n_S\to S$ remains finite type after any base change.\end{exercise}

\begin{exercise}\label{exr:vi26-09}Let $A$ be Noetherian. Show every localization $S^{-1}A$ is Noetherian.\end{exercise}

\begin{exercise}\label{exr:vi26-10}Let $A$ be Noetherian. Show every quotient $A/I$ is Noetherian.\end{exercise}

\begin{exercise}\label{exr:vi26-11}Use Hilbert's basis theorem to show $A[x_1,\dots,x_n]$ Noetherian when $A$ is.\end{exercise}

\begin{exercise}\label{exr:vi26-12}Show a finite-type algebra over a Noetherian ring is Noetherian.\end{exercise}

\begin{exercise}\label{exr:vi26-13}Show the ascending chain $(x_1)\subset(x_1,x_2)\subset\cdots$ in $k[x_1,x_2,\dots]$ is strict.\end{exercise}

\begin{exercise}\label{exr:vi26-14}Explain why the infinite disjoint union of points is not quasi-compact.\end{exercise}

\begin{exercise}\label{exr:vi26-15}Show a finite disjoint union of finite-type $k$-schemes is finite type over $k$.\end{exercise}

\begin{exercise}\label{exr:vi26-16}Show a fiber of a finite-type morphism is finite type over its residue field.\end{exercise}

\begin{exercise}\label{exr:vi26-17}Show local finite type is preserved by composition.\end{exercise}

\begin{exercise}\label{exr:vi26-18}Show finite type is preserved by composition.\end{exercise}

\begin{exercise}\label{exr:vi26-19}Show local finite type is preserved by base change.\end{exercise}

\begin{exercise}\label{exr:vi26-20}Show finite type is preserved by base change.\end{exercise}

\begin{exercise}\label{exr:vi26-21}If $X$ is locally Noetherian, show every open subscheme is locally Noetherian.\end{exercise}

\begin{exercise}\label{exr:vi26-22}If $X$ is Noetherian, show every open subscheme is quasi-compact.\end{exercise}

\begin{exercise}\label{exr:vi26-23}Show that a finite-type scheme over a field is Noetherian.\end{exercise}

\begin{exercise}\label{exr:vi26-24}Give a locally finite type but non-finite-type morphism different in notation from the chapter example.\end{exercise}

\begin{exercise}\label{exr:vi26-25}Give a finite-type morphism that is not finite as a morphism.\end{exercise}

\begin{exercise}\label{exr:vi26-26}Explain why finite type is a statement about generators, not about the number of points in a fiber.\end{exercise}
